\subsection{Leur vie à Ferrussac}

Jean et Thérèse Gaillard-Arnal de l’Oultre, à Pont Vieux, me donnèrent bien des
précisions sur l’installation de ces familles harkies. En Lozère, 160 familles environ
furent installées dans ce qu’on a appelé des « hameaux forestiers » : 52 à Villefort, 26
à Chadenet, 26 à Cassagnas-Haut, 26 à Culture, et enfin 26 à Meyrueis au lieu-dit La
Foulatière, un long « bancel », à 200m en amont de la chaussée de Roquedols au -
dessus de la route de Ferrussac, limité par le vallon de Roquedoulet.

Jean Gaillard, menuisier, vint à Meyrueis avec les camions apportant les baraques
préfabriquées et prêtes à être montées. C’est une entreprise de Saint-Pardoux-la-
Rivière en Dordogne, spécialisées dans la fabrication de villas, salles de classe,
baraques de chantier préfabriquées, qui obtint le marché.
Les familles harkies furent donc installées à Meyrueis le 15 mars 1963 dans ces
baraquements légers pour hébergement provisoire

\begin{quote}
\textit{« C’était, me dit-Jean Gaillard, des baraques légères à parois de bois isolées par
seulement 5cm de laine de verre ( nous sommes en 1963 ). Les familles disposaient de
2 ou 3 pièces suivant leur composition : une ou deux femmes -l’épouse et la
concubine- Chaque logement avait une cuisine avec évier et une cheminée. Il y avait
donc, en tout, 60 à 70 adultes accompagnés par le même nombre d’enfants : les
baraques se chauffaient rapidement mais se refroidissaient instantanément dès que le
feu « tombait », et l’été, par grand soleil il fallait arroser les toits pour faire baisser
la chaleur étouffante à l’intérieur.. Les toilettes étaient à l’extérieur. A côté des
logements une baraque servait de foyer où les hommes pouvaient se réunir après le
travail. Le soubassement subsiste encore »}
\end{quote}

Quant à Thérèse Gaillard, Melle. Arnal à l’époque, elle avait pendant sa scolarité, à la
fin des années 50, été très présente auprès d’une grand mère âgée et malade, l’aidant
pour les repas, la toilette etc. Elle venait d’obtenir un brevet d’aide-ménagère-
comptable. Elle fut embauchée pour travailler auprès des femmes harkies, d’abord
dans un hameau perdu au-dessus de Villefort, le Pouget, afin de faciliter leur
adaptation.

\begin{quote}
\textit{« J’ai posé ma candidature pour Meyrueis quand j’ai su qu’on allait implanter un
hameau près de Ferrussac. Elle a été retenue en partie parce que je disposais d’une
2cv indispensable pour les trajets. J’étais désormais « Monitrice d’Action Sociale ».
Lorsque la 2cv arrivait, toutes les femmes se disaient l’une à l’autre « C’est la
sociale ! ».
Une de mes plus grosses difficultés fut d’obtenir l’ouverture des fenêtres et des volets
pour aérer et laisser entrer le soleil et la lumière. Il y avait aussi les soins pour les
bébés : les femmes les emmaillotaient totalement, la tête seule dépassait, et les
nourrissons ne pouvaient bouger ni bras ni jambes. Car il y eut des naissances . Ce
ne fut pas une mince affaire d’emmener les femmes, avec leurs maris bien sûr, à la
maternité de Florac ou de Mende. Il fallait également leur apprendre à acheter aux
épiciers ambulants, des aliments et des légumes différents de ceux d’Algérie,
proposés sur place par Caiffa, « Pissou » de Florac, Rigal etc. Mme. Richard de
Ferrussac les voyait fréquemment en leur apportant le lait de ses vaches.
Il fallut leur apprendre à consulter le médecin ou à l’appeler, heureusement le Dr.
Juliette Cabanel était bien acceptée car c’était une femme. »}
\end{quote}
